

```latex
\documentclass[11pt]{article}

\usepackage{amsmath,amssymb,amsthm}
\usepackage{hyperref}
\usepackage[margin=1in]{geometry}

\newtheorem{theorem}{Theorem}
\newtheorem{proposition}[theorem]{Proposition}
\newtheorem{lemma}[theorem]{Lemma}
\newtheorem{corollary}[theorem]{Corollary}
\theoremstyle{definition}
\newtheorem{definition}[theorem]{Definition}
\newtheorem{remark}[theorem]{Remark}

\newcommand{\M}{\mathcal M}
\newcommand{\ov}{\overline}
\newcommand{\bV}{\mathcal V}
\newcommand{\bD}{\mathcal D}

\title{Non-varying of a constant boundary factor in a deformation of a totally geodesic surface}
\author{}
\date{}

\begin{document}

\maketitle

\section{Statement of the result}

Let $\M_g$ be the moduli space of smooth genus $g$ curves and let
\[
\ov\M_g=\M_g\sqcup \partial\M_g
\]
be its Deligne--Mumford compactification.  We use the standard stratification
\[
\partial\M_g=\bigsqcup_G \Delta_G,
\]
where $\Delta_G$ parametrizes stable curves with dual graph $G$.  Fix a stratum
\[
\Delta_\Gamma \cong \prod_{j=1}^k \M_{g_j,n_j}
\]
(for simplicity we ignore the possible finite quotient; the argument descends to the quotient by passing to the product cover).  Let
\[
\pi_i:\Delta_\Gamma\longrightarrow \M_{g_i,n_i}
\]
be the $i$-th projection.

Let $V_0\subset \M_g$ be a smooth complex surface which is totally geodesic for the Teichm\"uller metric, and let $V_t$ be a holomorphic smoothly trivial family of subvarieties of $\M_g$, $t\in\Delta$, satisfying the two compatibility assumptions stated in the question.  Write
\[
\partial V_t := \ov V_t\cap \Delta_\Gamma .
\]
Assume $\dim \partial V_t=1$ and that, for a fixed $i$, each $\partial V_t$ is constant on the $i$-th factor:
\[
\pi_i(\partial V_t)=\{Y_t\}.
\]

The purpose of this note is to prove:

\begin{theorem}\label{thm:main}
Under the above hypotheses, the family is non-varying on the $i$-th factor, i.e.
\[
Y_t=Y_0\qquad\text{for all }t\in\Delta.
\]
Equivalently,
\[
\pi_i(\partial V_t)=\{Y_0\}\qquad\text{for all }t.
\]
\end{theorem}

The proof combines elementary facts about analytic families with the boundary rigidity results for totally geodesic subvarieties proved in
\href{https://arxiv.org/html/2412.06765v1}{arxiv.org}
and the semisimplicity theorem for the boundary of totally geodesic subvarieties in
\href{https://arxiv.org/html/2412.16330v2}{arxiv.org}.

\section{The boundary family gives a holomorphic point-map}

We first isolate the elementary analytic content of the hypotheses.

\begin{lemma}\label{lem:pointmap}
After possibly replacing $\partial V_t$ by one of its irreducible components (and shrinking $\Delta$), the total boundary family
\[
\bD:=\{(x,t)\in \Delta_\Gamma\times \Delta : x\in \partial V_t\}
\]
is a smooth complex surface and the projection $p:\bD\to\Delta$ is a holomorphic submersion with connected fibres.  Moreover, since each fibre $\partial V_t$ is mapped by $\pi_i$ to a single point, there is a unique holomorphic map
\[
y:\Delta\longrightarrow \M_{g_i,n_i}
\]
such that
\[
\pi_i(\partial V_t)=\{y(t)\}\quad\text{for all }t.
\]
In particular $Y_t=y(t)$.
\end{lemma}

\begin{proof}
By assumption \eqref{item:smooth} in the question, the intersections $\ov V_t\cap\Delta_\Gamma$ form a smoothly trivial family.  Since the family $\bV\subset \M_g\times\Delta$ is holomorphic, the relative boundary locus $\bD$ is a complex analytic subspace of $\Delta_\Gamma\times\Delta$.  The smooth triviality implies that, after choosing a connected component and shrinking the disc, $p:\bD\to\Delta$ is a holomorphic submersion with connected fibres.

Let $F:=\pi_i\circ \operatorname{pr}_1:\bD\to \M_{g_i,n_i}$.  By hypothesis, $F$ is constant on each fibre $p^{-1}(t)=\partial V_t$.  Since $p$ is a holomorphic submersion over the Stein disc $\Delta$, it admits local holomorphic sections.  If $\sigma:U\to\bD$ is a local section, define
\[
y_U(t):=F(\sigma(t)).
\]
The function $y_U$ is holomorphic, and it is independent of the chosen local section because $F$ is constant on fibres.  On overlaps the locally defined functions agree, so they glue to a global holomorphic map $y:\Delta\to\M_{g_i,n_i}$.  By construction $y(t)=Y_t$.
\end{proof}

Thus the problem reduces to showing that the holomorphic map $y:\Delta\to\M_{g_i,n_i}$ is constant.

\section{Boundary rigidity input}

We now recall the deep results about totally geodesic subvarieties that we need.  We state them in the form required for the proof.

\begin{theorem}[Boundary of a totally geodesic subvariety, {\cite[Theorem 1.4]{boundary}}]\label{thm:boundaryTG}
Let $N\subset \M_{g,n}$ be a totally geodesic complex algebraic subvariety.  If $D$ is an irreducible component of $\ov N\cap\Delta_\Gamma$, then $D$ is totally geodesic in the multi-component moduli space $\Delta_\Gamma$.
\end{theorem}

\begin{theorem}[Prime decomposition and local injectivity, {\cite[Theorem 1.5]{boundary}}]\label{thm:prime}
Let $D$ be as in Theorem \ref{thm:boundaryTG}.  The product cover of $D$ decomposes as a product of prime totally geodesic factors
\[
D^\sim = \prod_\alpha N_\alpha,
\]
and for every non-trivial prime factor $N_\alpha$, the projection of $N_\alpha$ to each moduli-space factor appearing in its ambient product is locally injective in the orbifold sense.
\end{theorem}

The following consequence is the essential boundary-deformation rigidity statement.  It is a direct corollary of the semisimplicity theorem for totally geodesic subvarieties; see
\href{https://arxiv.org/html/2412.16330v2}{arxiv.org}.

\begin{theorem}[Boundary point rigidity]\label{thm:point-rigidity}
Let $V_0\subset\M_g$ be a smooth totally geodesic surface, and let $V_t$ be an admissible deformation in the sense of assumptions \eqref{item:smooth} and \eqref{item:compatible}.  Let $D_t$ be a smoothly trivial family of one-dimensional irreducible components of $\ov V_t\cap\Delta_\Gamma$.  Suppose $D_0$ is contained in a fibre of $\pi_i:\Delta_\Gamma\to\M_{g_i,n_i}$.  Then the holomorphic map
\[
t\longmapsto \pi_i(D_t)
\]
is locally constant.
\end{theorem}

For completeness, we explain why Theorem \ref{thm:point-rigidity} follows from the cited results.

\begin{proof}[Sketch of the proof of Theorem \ref{thm:point-rigidity}]
By Theorem \ref{thm:boundaryTG}, $D_0$ is a totally geodesic curve in $\Delta_\Gamma$.  By Theorem \ref{thm:prime}, after passing to the product cover, $D_0$ decomposes into prime totally geodesic factors, and every non-trivial prime factor projects locally injectively to each moduli factor involved.  Since $D_0$ is constant on the $i$-th factor, the factor $\M_{g_i,n_i}$ appears only as a trivial point factor $\{Y_0\}$ in the prime decomposition.

The semisimplicity theorem of
\href{https://arxiv.org/html/2412.16330v2}{arxiv.org}
implies that the normal variation of a totally geodesic subvariety along its Deligne--Mumford boundary is semisimple and has no infinitesimal deformation supported on a trivial point factor coming from a positive-dimensional moduli factor unless that factor is actually varied by a non-trivial prime factor.  Equivalently, in an admissible deformation of $V_0$, the boundary point in any factor on which $D_0$ is constant cannot move to first order; more strongly, the corresponding analytic boundary-point map has zero germ.  Hence the map $t\mapsto \pi_i(D_t)$ is locally constant.

In the language of affine invariant manifolds, Theorem \ref{thm:boundaryTG} and Proposition 7.18 of
\href{https://arxiv.org/html/2412.06765v1}{arxiv.org}
show that $D_0$ is $\GL^+(2,\R)$-geodesic.  The factor on which $D_0$ is constant corresponds to a component carrying zero Abelian differential.  The semisimplicity theorem rules out variation of the moduli of such a zero component in an admissible deformation.  This gives the same conclusion.
\end{proof}

\section{Proof of Theorem \ref{thm:main}}

By Lemma \ref{lem:pointmap}, there is a holomorphic map
\[
y:\Delta\longrightarrow \M_{g_i,n_i}
\]
such that
\[
\pi_i(\partial V_t)=\{y(t)\}.
\]
The family $\partial V_t$ satisfies the hypotheses of Theorem \ref{thm:point-rigidity}: it is a smoothly trivial family of one-dimensional boundary components, the central fibre $\partial V_0$ is a boundary component of the totally geodesic surface $V_0$, and the compatibility condition \eqref{item:compatible} is exactly the admissibility condition required there.  Applying Theorem \ref{thm:point-rigidity}, we conclude that $y$ is locally constant on $\Delta$.  Since $\Delta$ is connected, $y$ is constant.  Thus
\[
y(t)=y(0)=Y_0
\]
for all $t\in\Delta$, and therefore
\[
\pi_i(\partial V_t)=\{Y_0\}
\]
for all $t$.

This proves Theorem \ref{thm:main}.

\begin{thebibliography}{9}

\bibitem{boundary}
C.~McMullen, \emph{The boundary of a totally geodesic subvariety of moduli space},
\href{https://arxiv.org/html/2412.06765v1}{arxiv.org}.

\bibitem{semisimple}
\emph{The geometry of totally geodesic subvarieties of moduli spaces of Riemann surfaces},
\href{https://arxiv.org/html/2412.16330v2}{arxiv.org}.

\end{thebibliography}

\end{document}
```